%% 17_cdac.tex -- Chapter 17: Capacitive DACs and References
%% Status: OUTLINE (build stage 4). Replace the outline with the full chapter when it is written.
\chapter{Capacitive DACs and References}\label{ch:cdac}
\chapterlead{The capacitive DAC is both the sample-and-hold and the reference ladder of a SAR converter; its matching sets the linearity.}

\begin{specnote}[Chapter status: outline]
This chapter is an outline. It will be written in build stage 4 (the rest of Part III with Labs 7, 8 and 10). The topics, the design question and the labs below are fixed; the derivations, worked examples and figures are still to come.
\end{specnote}

\begin{objectives}
\begin{itemize}
  \item compare binary-weighted, split (bridge) and C-2C capacitor arrays;
  \item compute the switching energy of the conventional, monotonic, $V_\mathrm{cm}$-based and merged-capacitor schemes;
  \item size the unit capacitor from matching and from $kT/C$;
  \item choose between MOM and MIM capacitors in \sgx{} and lay them out in common centroid;
  \item design reference buffering and decoupling.
\end{itemize}
\end{objectives}

\begin{insight}[Design question]
How small can the unit capacitor be for 12-bit linearity, and which switching scheme minimises the DAC energy?
\end{insight}

\section{Capacitor array topologies}\label{sec:ch17-1}
Binary arrays, split arrays with a bridge capacitor, and C-2C ladders: area, linearity and parasitic sensitivity.

\section{Switching schemes and energy}\label{sec:ch17-2}
Conventional, monotonic, $V_\mathrm{cm}$-based and merged-capacitor switching; the average energy per conversion of each.

\section{Unit-capacitor sizing}\label{sec:ch17-3}
Matching-limited versus $kT/C$-limited size; why matching usually wins for SAR converters above 10 bits.

\section{Capacitors in \texorpdfstring{\sgx{}}{SG13G2}}\label{sec:ch17-4}
MOM and MIM options, their density and matching, and common-centroid layout of the array.

\section{References}\label{sec:ch17-5}
The DAC draws charge from the reference every bit cycle; on-chip buffers versus large decoupling capacitors.

\begin{tryit}[Labs for this chapter]
Lab~8 (\cref{sec:lab8}) builds a CDAC with monotonic switching; Lab~14 (\cref{sec:lab14}) lays it out.
\end{tryit}

\paragraph{Read with this chapter} Pelgrom~\cite{pelgrom}, Moser and Pretl~\cite{jkusar} (an open non-binary 12-bit CDAC).
